\documentclass[11pt,a4paper]{article}
\usepackage[italian]{babel}
\usepackage[margin=2.5cm]{geometry}
\usepackage{parskip}
\usepackage{amsmath, amssymb, amsthm}
\usepackage{booktabs}
\usepackage{array}
\usepackage{xcolor}
\usepackage{needspace}
\usepackage{enumitem}
\usepackage{listings}
\usepackage{tikz}
\usepackage[most]{tcolorbox}
\usepackage[hidelinks]{hyperref}
\usetikzlibrary{decorations.pathreplacing, shapes.geometric, arrows.meta, positioning}

% Niente vedove/orfane: i paragrafi non lasciano righe isolate a fine/inizio pagina
\widowpenalty=10000
\clubpenalty=10000

% Elenchi più compatti
\setlist{itemsep=2pt, parsep=0pt, topsep=4pt, partopsep=0pt}

% --- Riquadri con bordo nero, senza numero (si spezzano tra due pagine) ---
% Uso: \begin{definizione} ... \end{definizione}
%      \begin{teorema}[Nome] ... \end{teorema}  (il nome è facoltativo)
\tcbset{nota/.style={enhanced jigsaw, breakable, colback=white,
  colframe=black, boxrule=0.5pt, sharp corners}}
\NewTColorBox{definizione}{}{nota,
  before upper={\textbf{Definizione.}\ }}
\NewTColorBox{teorema}{o}{nota,
  before upper={\textbf{Teorema\IfValueT{#1}{ (#1)}.}\ }}
\NewTColorBox{proposizione}{o}{nota,
  before upper={\textbf{Proposizione\IfValueT{#1}{ (#1)}.}\ }}
\NewTColorBox{proprieta}{o}{nota,
  before upper={\textbf{Proprietà\IfValueT{#1}{ (#1)}.}\ }}
\NewTColorBox{assioma}{o}{nota, boxrule=1pt,
  before upper={\textbf{Assioma\IfValueT{#1}{ (#1)}.}\ }}

% --- Ambienti semplici (senza riquadro) ---
% Il titolo sta in cima al blocco, anche se il contenuto inizia con un riquadro o
% con codice; e non resta mai da solo in fondo a una pagina.
% Uso: \begin{esempio}[Titolo facoltativo] ... \end{esempio}
\newcommand{\newrunin}[2]{%
  \NewDocumentEnvironment{#1}{o}{%
    \par\Needspace{4\baselineskip}\addvspace{\medskipamount}\noindent
    \textbf{#2}\IfValueT{##1}{ (##1)}\textbf{.}\ \ignorespaces}%
  {\par\addvspace{\medskipamount}}}
\newrunin{esempio}{Esempio}
\newrunin{esercizio}{Esercizio}
\NewTColorBox{osservazione}{}{nota,
  before upper={\textbf{Osservazione.}\ }}

% V e F nelle tavole di verità
\newcommand{\V}{\textsf{V}}
\newcommand{\F}{\textsf{F}}

% --- Codice C ---
% Uso: \begin{codice} ... \end{codice}      codice C numerato
%      \begin{comando} ... \end{comando}    comandi da terminale
%      \begin{risultato} ... \end{risultato} output di un programma
%      \lstinline|printf("ciao");|           codice dentro al testo
\lstdefinestyle{c}{
  language=C,
  basicstyle=\ttfamily\small,
  keywordstyle=\color{blue}\bfseries,
  commentstyle=\color{gray}\itshape,
  stringstyle=\color{red!70!black},
  numbers=left,
  numberstyle=\tiny\color{gray},
  numbersep=8pt,
  columns=flexible,
  keepspaces=true,
  breaklines=true,
  showstringspaces=false,
  tabsize=4,
  morekeywords={include, define},
  % lettere accentate nei commenti
  literate={à}{{\`a}}1 {è}{{\`e}}1 {é}{{\'e}}1 {ì}{{\`i}}1
           {ò}{{\`o}}1 {ù}{{\`u}}1 {È}{{\`E}}1 {À}{{\`A}}1
}
\lstset{style=c}
\newtcblisting{codice}{listing only, nota, left=6mm,
  listing options={style=c}}
\newtcblisting{comando}{listing only, nota,
  listing options={basicstyle=\ttfamily\small, numbers=none, language={},
    columns=flexible, keepspaces=true, breaklines=true}}
\newtcblisting{risultato}{listing only, enhanced jigsaw, breakable,
  colback=black!5, colframe=black, boxrule=0.5pt, sharp corners,
  listing options={basicstyle=\ttfamily\small, numbers=none, language={},
    columns=flexible, keepspaces=true, breaklines=true}}

% --- Diagrammi di flusso (simboli standard) ---
\tikzset{
  terminale/.style = {ellipse, draw, minimum width=2.5cm, minimum height=0.9cm},
  io/.style        = {trapezium, trapezium left angle=70, trapezium right angle=110,
                      draw, minimum width=2.5cm, minimum height=0.9cm},
  processo/.style  = {rectangle, draw, minimum width=2.5cm, minimum height=0.9cm},
  decisione/.style = {diamond, draw, aspect=2, minimum width=2.5cm},
  freccia/.style   = {thick, -{Stealth}}
}

\title{Costanti letterali}
\author{Marco Cerbella}
\date{Lezione 9 -- 12 ottobre 2026}

\begin{document}
\maketitle

\section{Costanti letterali}

\begin{definizione}
Un \emph{literal} (costante letterale) è un token che denota un valore fisso, che può essere un intero, un numero in virgola mobile, un carattere o una stringa. Il tipo di un literal è determinato dal suo valore e dalla sua notazione.
\end{definizione}

\section{Costanti intere}

Una costante intera si può scrivere come un normale numero decimale oppure in notazione ottale o esadecimale. La notazione si specifica con un \emph{prefisso}.

\begin{itemize}
    \item \textbf{Decimale}: inizia con una cifra diversa da zero.
    \item \textbf{Ottale} (base 8, cifre $0, \dots, 7$): inizia con uno zero iniziale. Il valore è $c_n 8^n + c_{n-1} 8^{n-1} + \dots + c_1 8^1 + c_0 8^0$.
    \item \textbf{Esadecimale}: inizia con il prefisso \lstinline|0x| o \lstinline|0X|. Le cifre da A a F possono essere maiuscole o minuscole.
\end{itemize}

\begin{esempio}
\leavevmode
\begin{itemize}
    \item $047 = 4 \cdot 8^1 + 7 \cdot 8^0 = 39$.
    \item \lstinline|0xff|, \lstinline|0Xff|, \lstinline|0xFF| e \lstinline|0XFF| valgono tutte $15 \cdot 16^1 + 15 \cdot 16^0 = 255$.
\end{itemize}
\end{esempio}

\subsection{Tipo delle costanti intere}

Il tipo di una costante è determinato nel momento in cui ne viene definito il valore. Le costanti intere degli esempi precedenti hanno di solito tipo \lstinline|int|. Se il valore è troppo grande, il compilatore assegna il primo tipo, in una gerarchia, abbastanza grande da rappresentarlo. Ad esempio, la costante decimale \lstinline|50000| ha tipo \lstinline|int|, perché il massimo valore possibile per uno \lstinline|short| (2 byte) è $32\,767 = 2^{15} - 1$.

Si può influenzare esplicitamente il tipo di una costante con un \emph{suffisso}:
\begin{itemize}
    \item \lstinline|512U|: \lstinline|unsigned int|;
    \item \lstinline|0Xf0fUL|: \lstinline|unsigned long|;
    \item \lstinline|0777ll|: \lstinline|long long|;
    \item \lstinline|123uLL|: \lstinline|unsigned long long|.
\end{itemize}
Maiuscole e minuscole sono equivalenti (\lstinline|ll| e \lstinline|LL| sono la stessa cosa).

\begin{codice}
int a = 512U;
int b = 1LL;
int c = 7UL;
\end{codice}

\section{Costanti in virgola mobile}

Le costanti in virgola mobile si scrivono in notazione decimale o esadecimale. Una costante è una sequenza di cifre decimali contenente un punto decimale. Si può moltiplicare il valore per una potenza di $10$, come nella notazione scientifica: la potenza di $10$ si indica con un esponente introdotto dalla lettera \lstinline|e| o \lstinline|E|.

\begin{esempio}
\lstinline|2.34E5| vale $2.34 \cdot 10^5$ e \lstinline|67e-12| vale $67.0 \cdot 10^{-12}$.
\end{esempio}

Il punto decimale può anche essere il primo o l'ultimo carattere: \lstinline|10.| e \lstinline|.234E6| sono numeri validi. Il numero \lstinline|10| senza punto decimale è invece una costante \emph{intera}, non in virgola mobile.

\subsection{Forzare il tipo}

Il tipo di default di una costante in virgola mobile è \lstinline|double|. Con il suffisso \lstinline|F| o \lstinline|f| si assegna il tipo \lstinline|float|; con il suffisso \lstinline|L| o \lstinline|l| il tipo \lstinline|long double|.

\section{Costanti carattere}

Una costante carattere è un carattere racchiuso tra apici singoli: \lstinline|'a'|, \lstinline|'0'|, \lstinline|'*'|. Per rappresentare l'apice, il backslash e il ritorno a capo si usano \lstinline|'\''|, \lstinline|'\\'| e \lstinline|'\n'|. Le costanti carattere hanno tipo \lstinline|int|.

\subsection{Sequenze di escape}

\begin{center}
\begin{tabular}{@{}lll@{}}
\toprule
Sequenza & Valore del carattere & Azione sul dispositivo di output \\
\midrule
\verb|\'| & apice singolo (\lstinline|'|) & stampa il carattere \\
\verb|\"| & virgolette (\lstinline|"|) & stampa il carattere \\
\verb|\?| & punto interrogativo & stampa il carattere \\
\verb|\\| & backslash & stampa il carattere \\
\verb|\a| & alert & produce un segnale acustico o visivo \\
\verb|\b| & backspace & riporta la posizione attiva indietro di un carattere \\
\verb|\f| & form feed & porta la posizione all'inizio della pagina successiva \\
\verb|\n| & line feed & porta la posizione all'inizio della riga successiva \\
\verb|\r| & carriage return & porta la posizione all'inizio della riga corrente \\
\bottomrule
\end{tabular}
\end{center}

\section{Stringhe letterali}

Una \emph{stringa letterale} è una sequenza di caratteri (e/o sequenze di escape) racchiusa tra virgolette doppie, ad esempio \lstinline|"Hello world!\n"|. Come le costanti carattere, può contenere tutti i caratteri del set di caratteri sorgente.

Una stringa letterale è un \emph{array statico di char} che contiene i codici dei caratteri seguiti dal terminatore di stringa, il carattere nullo \lstinline|\0|. La stringa vuota \lstinline|""| occupa esattamente un byte di memoria, che contiene il carattere nullo di terminazione.

\begin{codice}
char helloWorld[128] = "Hello World!\n";
printf("Print the string: %s\n", helloWorld);
\end{codice}

\section*{Riferimenti}

Sul libro: sezione 15.6 (pagina 339) e pagina 46, sezione 9.10.

\end{document}
